
\section{Pseudo-quadratic-like maps and pseudo-puzzle}\label{pseudo-puzzle}

\subsection{Pseudo-quadratic-like maps}

For a more general and detailed discussion of $\psi$-ql maps, see \cite{K}.

Suppose that $\BU'$, $\BU$ are disks,
and $i:  \BU' \to \BU$ is a holomorphic immersion,
and $f: \BU' \to \BU$ is a degree $d$ holomorphic branched cover.
Suppose further that there exist full continua $K \Subset \BU$ and $K'\Subset \BU'$
such that $K' = i^{-1}(K) = f^{-1}(K)$.
Then we say that $F=(i,f):  \BU'\ra (\BU, \BU)$  is a
$\psi$-{\it quadratic-like}
 ($\psi$-{\it ql}) map  with filled Julia set $K$.

\begin{lem}[\cite{K}]\label{ql extension}
Let $F=(i,f)\from\BU'\ra \BU$ be a $\psi$-ql map of degree $d$
 with filled Julia set $K$.
Then $i$ is an embedding in  a neighborhood of $K' \equiv f^{-1}(K)$,
and the map  $g\equiv  f\circ i^{-1}\from U'\ra U$ near $K$ is quadratic-like.

\ssk Moreover, the domains $U$ and $U'$ can be selected in such a way that
$\mod(U\sm i(U'))\geq \mu (\mod(\BU\sm K)) >0$.
\end{lem}



There is a natural $\psi$-ql map $\BU^n\ra \BU^{n-1}$,
the ``restriction'' of $(i,f)$ to $\BU^n$.
Somewhat loosely, we will use the same notation $F=(i,f)$ for this restriction.

 Let us normalize  the $\psi$-quadratic-like  maps under consideration
so that $\diam K'= \diam K=1$, both $K$ and $K'$ contain 0 and 1,
$0$ is the critical point of $f$, and $i(0)=0$.
Let us endow the space of $\psi$-quadratic-like maps (considered up to independent rescalings in the domain and the range)
with the Carath\'eodory topology.
In this topology, a sequence of normalized maps $(i_n, f_n): \BU_n'\ra\BU_n$ converges to $(i,f): \BU'\ra \BU$ if
the pointed domains $(\BU_n', 0)$ and $(\BU_n, 0)$ converge to $\BU'$ and $\BU$ respectively,
and the maps $i_n$, $f_n$ converge respectively to $i$, $f$,
uniformly on compact subsets of  $\BU'$.


\begin{lem}[compare \cite{McM}]\label{compactness}
  Let $\mu>0$. Then the space of $\psi$-PL maps $F=(i,f): \BU'\ra \BU$ such that the  Julia set $K$
is connected and  $\mod(\BU \sm K)\geq \mu$ is compact.
\end{lem}

\begin{proof}
 Let  $X_n= i^{-1}\{0,1\}$,  $Y_n= f^{-1}\{0,1\}$.
Note that both sets consist of at most 2 points and are contained in $\D$.

Then we can select a subsequence of domains $\BU_n'$, $\BU_n$ Carath\'eodory converging to some
domains $\BU'$, $\BU$, while the sets $X_n$ and $Y_n$ converge in the Hausdorff metric to some sets
$X\subset \BU'$  and $Y\subset \BU'$ that consist of at most two points and are contained in $\D$
 (we will keep the same notation for the subsequence).
Since  the maps $i_n| \BU_n'\sm X_n$ and $f_n| \BU_n'\sm Y_n$ do not assume values $0$ and $1$,
they form normal families on $\BU'\sm X $. Since these families are bounded on the sets $K_n$,
they are uniformly bounded on compact sets of $\BU'\sm (X\cup Y)$.
By the Maximum Principle, they are normal on the whole domain $\BU'$.


Let $i$ and $f$ be some limit functions of the sequences $i_n$ and $f_n$.
These functions are non-constant since they assume values $0$ and $1$.
Then $i$ is an immersion as a non-constant limit of immersions.
Also, $f: \BU'\ra \BU$ is a branched covering of degree at most 2.
% Let us also pass to a subsequence to make the Julia sets $K_n'$  converge in the Hausdorff metric to some set $K'$.
Moreover,  $K'\Subset \BU$ since $\mod(\BU'\sm K')\geq \mu/2$. Hence $0\in \BU'$, and it is a critical point of $f$.
It follows that $\deg f=2$, and we are done.
\end{proof}


\subsection{Pseudo-puzzle}\label{psi-puzzle}

\subsubsection{Definitions}
   Let $(i,f) : \BU'\ra \BU$ be a $\psi$-ql map.
By Lemma~\ref{ql extension}, it admits a quadratic-like restriction $U'\ra U$ to a neighborhood of
its (filled) Julia set  $K=K_\BU$.
Here $U'$ is embedded to $U$, so we can identify $U'$ with $i(U')$
and $ f: U'\ra U $ with  $f \circ i^{-1} $.

Assume that $K$ is connected and both fixed points of $f$ are repelling.
Then we can cut $U$ by external rays landing at the $\alpha$-fixed point and consider the corresponding Yoccoz puzzle.
%
%  Assume $f$ is renormalizable with the first pre-renormalization $g$.
% As in \S \ref{can renorm},   let $K_i(g)$, $i=0,1,\dots, p-1$,
% stand for  the cycle of its little Julia sets, and let $\tl K(g)=\cup K_i(g)$.
%
%In fact, besides {\it dynamical} puzzle pieces that appear in the puzzle,
%we can consider {\it geometric} puzzle pieces defined as
%unions of some dynamical pieces of the same depth.
%%%topological disks bounded by finitely many rational rays and equipotentials.
%In what follows we will use the term ``puzzle piece'' in this more general sense.

Given a (geometric) puzzle piece $Y$ of depth $m$,
recall that $K_Y$ stands for $Y\cap K(f)$ and  $S_Y = \cl (K(f) \sm K_Y)$.
Let  $\YY$ stand for the space of paths $\de: [0,1] \to \BU^m \sm S_Y$
such that:
\begin{itemize}
  \item $\de(0)\in Y$,
  \item if $\de(t)\in K_Y$, then the restriction $\de|\, [0,t]$ is homotopic rel endpoints
       to a path contained in $Y$.\footnote{This condition can be replaced with a more restrictive one:
 After the first exit from $Y$, the path never intersect the Julia set $K(f)$
  (though it is allowed to return back to $Y$).}
\end{itemize}
Let $\BY$ be the space of paths $\de\in \YY$ modulo  homotopy through $\YY$ with $\de(1)$ fixed.
Define the projection $\pi_Y: \BY\ra \BU^m$ by  $[\de]\mapsto \de(1)$.
% In the same way as in \S \ref{can renorm},
One can see that $\BY$ is a Riemann surface,
and $\pi_Y$ is an immersion such that
   $Y$ lifts to  a disk $\hat Y\subset \BY$ which is  homeomorphically projected onto $Y$.
 Thus, we can identify $\hat Y$ with $Y$; in particular,  $K_Y$ is embedded into $\BY$.
%  \item $\pi_Y: \BY \sm \hat Y  \ra  V\sm (Y\cup Q_Y)$ is a covering.

   The Riemann surface $\BY$ will be called the {\it pseudo-piece} (``$\psi$-piece'') associated with $Y$.

\ssk The $\psi$-pieces can also be defined in a different way.
Let us consider the topological annulus $A=\BU^m  \sm K(f)$ and its universal covering $\hat A$.
Let $Y_i$ be the components of $Y\sm K_Y$. There are finitely many of them,
and each $Y_i$ is simply connected. Hence they can be embedded into $\hat A$.
Select such an embedding $e_i: Y_i\ra \hat A_i$ where $\hat A_i$ stands for a copy of $\hat A$.
Then the $\psi$-piece is obtained by gluing the $A_i$ to $Y$ by means of $e_i$, i.e.,
$\BY= Y \sqcup_{e_i} \hat A_i.$

\begin{lem}
  The above two definitions of $\psi$-pieces are equivalent.
\end{lem}

\begin{proof}
  Let $\BY$ be a $\psi$-piece according to the first definition.
The puzzle piece $Y$ is embedded into $\BY$ by
associating to a point $y \in Y$ the constant path $\de(t)\equiv y$.

Let us realize the universal covering $\hat A_i\ra A$ as the space of paths
in $A$ that begin in $Y_i$ rel homotopy through such paths fixing the terminal endpoint.
(This realization is legitimate since $Y_i$ is simply connected.)
This provides us with an embedding $\phi_i: \hat A_i\ra \BY $

The embeddings $\phi_i$ have disjoint images.
Indeed, all points of $\di Y\cap K(f)$ are dividing
and thus belong to $S_Y$. Hence, if we take two paths $\de_1: [0,1]\ra \BU^m\sm K(F)$
and $\de_2: [0,1]\ra  \BU^m\sm K(f)  $ as above
representing points in $\hat A_i$ and $\hat A_j$ ($i\not= j$) with a common endpoint,
then they ``surround'' some piece of $S_Y$, and hence represent different points in $\BY$.

Moreover,  the image $\phi_i(\hat A_i)$ overlaps with  $Y$ by $Y_i$.
 Hence we obtain an embedding of $Y\sqcup_{e_i} \hat A_i$ into $\BY$.

Let us show that this embedding is surjective.
Take  a path $\de\in \tilde\YY$ representing some point of $\BY$,
and let $\tau\in [0,1]$ be the last parameter for which $\de(\tau)\in K_Y$.
Since the path $\de: [0,\tau]\ra \BU^m$ is trivial
(i.e., it can be pulled to $Y$ in $\BU^m\sm S_Y$ rel endpoints),
the restriction $\de: [\tau, 1]\ra \BU^m \sm S_Y$ (appropriately reparametrized)
represents the same point in $\BY$ as the original path.
Moreover, if $\tau\not=1$, we can replace it with an equivalent path
%pull $\de(\tau)$ off the Julia set and make the path
$\de: [\tau+\eps , 1]\ra \BU^m \sm S_Y$
which is disjoint from the Julia set $K(f)$.
As the latter path represents a point in some $\hat A_i$, we are done.
\end{proof}



\sss{Naturality}

% The following properties follow immediately from the definitions:

\begin{lem}\label{psi-maps}
\begin{itemize}
\item [(i)]  Consider two puzzle pieces $Y$ and $Z$ such that the map $f: Y\ra Z$ is a branched covering
of degree $k$ (where $k=1$ or $k=2$ depending on whether $Y$ is off-critical or not).
% Then  the inverse map $f^{-1}: Z\ra Y$ extends to a $k$-valued immersion
Then there exists an induced map $\f: \BY \ra \BZ$
which  is a branched covering of the same degree $k$.


\item [(ii)] Given two puzzle pieces $Y\subset Z$, the inclusion $i: Y\ra Z$ extends to an immersion $\i: \BY\ra \BZ$.
\end{itemize}
\end{lem}

\begin{proof}
  Both properties follow easily from either definition of the $\psi$-pieces.
Let us, for instance, use the second definition.

(i)
Let $\depth Y=m$, $\depth Z=m-1$. Let us consider the degree $k$ branched covering
$$
    f: (\BU^m, Y, K_Y) \ra (\BU^{m-1}, Z, K_Z).
$$
The components $Y_i$ of $Y\sm K_Y$
are univalently mapped onto components $Z_{j(i)}$ of $Z\sm K_Z$, where the map
$j=j(i)$ is $k$-to-1.
This map extends to an isomorphism map $\hat A_i\ra \hat B_j $
of the corresponding universal coverings, which glue together into a branched covering
$\BY\ra \BZ$ of degree $k$.

(ii) Let $\depth Y=m$, $\depth Z=n$.
Let us  consider the immersion
$$
  i: (\BU^m, Y, K_Y) \ra (\BU^n, Z, K_Z).
$$
The components $Y_i$ are embedded by $i$ into some components $Z_{j(i)}$,
where the map $j=j(i)$ is surjective but not necessarily injective.
These embeddings extend to immersions $\hat A_i\ra \hat B_j$ that
glue together into an immersion $\BY\ra \BZ$.
\end{proof}

\sss{Moduli}

 Given two puzzle pieces $Z\Subset Y$, we let
$$
  \psimod (Y,Z) = \mod(\BY\sm K_Z).
$$

Lemma \ref{psi-maps} implies:

\begin{lem}\label{psi-mod}
\begin{itemize}
  \item [(i)]  Consider two pairs of puzzle pieces $(Y',Y)$ and $(Z',Z)$
such that the map $f: (Y',Y)\ra (Z',Z)$ is a branched covering
of degree $k$ (on both domains).
Then $$\psimod(Z',Z)=k\, \psimod(Y',Y).$$

\item [(ii)] Given a nest of three puzzle pieces $W\subset Z\subset  Y$,
we have $$\psimod(Z,W)\leq \psimod(Y,W).$$
\end{itemize}
\end{lem}


\subsubsection{Boundary of puzzle pieces}\label{boundary}
Let us mention in conclusion, that the ideal boundary of a pseudo-puzzle $\BY$
is tiled by (finitely many) arcs $\lambda_i\subset \di \hat A_i$
that cover the ideal boundary of $\BU^m$ (where $m=\depth Y$)
and arcs $\xi_i, \eta_i \subset \di \hat A_i$ mapped onto the Julia set $J(f)$.
The arc $\la_i$ meets each $\xi_i,\, \eta_i$ at a single boundary point
corresponding to a path $\de: [0,1)\mapsto A$ that wraps around $K(f)$ infinitely many times,
while $\eta_i$ meets $\xi_{i+1}$ at a vertex $v_i\in Y\cap K(f)$.
We say  that the arcs $\la_i$ form the {\it outer boundary } (or ``$O$-boundary'')
 $\di_O\BY$ of the puzzle piece $\BY$,
while the arcs $\xi_i$ and $\eta_i$
form its {\it $J$-boundary} $\di_J \BY$.
Given a vertex $v=v_i$ of a puzzle piece $Y$, let $\di^v\BY=\eta_i\cup \xi_{i+1}$
stand for the part of the $J$-boundary of $\BY$ attached to $v$.
%We let $\hat S_Y^v$ be the part of the $J$-boundary
%attached to the vertex $v$ of $Y$.

Note that the immersion
constructed in Lemma \ref{psi-maps} extends continuously to the boundary of the
puzzle piece $\BY$.
%  minus finitely many points where the $J$-boundary meets the $O$-boundery.
(However, $\i(\di \BY)$ is not contained in $\di\BZ$, unless $Z=Y$.)
In what follows we will assume this extension without further comment.


% In what follows it will be convenient to consider possibly disconnected curves.
A {\it multicurve} in some space $X$ is a continuous map $\gamma: \cup_{k=1}^l [s_k, t_k]\ra X$
parametrized by a finite union of disjoint intervals $[s_k, t_k]\subset \R$.%
\footnote{We allow that the boundary points of a  multicurve in a pseudo puzzle $\BY$ belong to  $\di_J\BY$.}
Note that  multicurves are ordered.
A multicurve in a puzzle piece $\BY$  is called {\it horizontal} if
$$
   \gamma(s_1)\in \di^{v_0} \BY, \quad  \gamma(t_k), \gamma(s_{k+1})\in \di^{v_k}\BY,\ k=1, \dots, l-1,
                  \quad \gamma(t_l)\in \di^{v_l} \BY
$$
for some vertices $v_k$ of $\BY$, $k=0, \dots, l$. We say that such a multicurve ``connects''
$\di^{v_0} \BY $ to $\di^{v_l} \BY$.
The following statement motivates introduction of  multicurves:

\begin{lem}\label{lifts}
Let $v$ and $w$ be two vertices of a geometric puzzle piece $Y\in\YY^m(l)$.
Then any curve $\gamma$ in $\BU^l$ connecting $S_Y^v$ to $S_Y^w$
contains a multicurve $\gamma'$  that lifts to a multicurve  $\gamma^*$ in $\BY$
connecting $\di_v\BY$ to $\di_w\BY$.
\end{lem}



Given two vertices $v$ and $w$ of $Y$, let $\GG_Y(v,w)$ stand for the family of horizontal multicurves in $\BY$
connecting $\di^v(\BY)$ to $\di^w(\BY)$.
Finally, let
$$
         \Bd_Y(v,w)= \LL(\GG_Y(v,w))
$$
stand for the extremal distance between the corresponding parts of $J$-boundary of $\BY$.
% Then we immediately conclude:

\begin{lem}\label{Bd}
  If  $f|Y$ is univalent,
then $\Bd_Y(v,w)=\Bd_{f(Y)}(fv, fw). $
\end{lem}


\comm{****************

\subsection{Pseudo-summary}\label{psi-summary}

   We can now recast the Summary of \S \ref{summary} in the pseudo-context.
We start with a $u\psi$-PL map $\bf$,
restrict it to a unicritical PL map $f$ and construct the associated dynasty of kingdoms.
We fix an arbitrary $m$, assume that the height of the dynasty is greater than $\log_2 m +5$,
and take some $n>\log_2 m +5$.
Then for any domain $\La =\La_k$, we obtain  the map $\Psi=\Psi_k$ \ref{3 domains}
satisfying properties (P1)-(P4). Applying to it the $\psi$-functor, we obtain:

\begin{itemize}

\item [(P1$'$)] An even-valued immersion ${\Bpsi}^{-1}: \BV^0 \ra \BUps$
 of degree bounded by $d^{2n+m}$;

\item[(P2$'$)] An even-valued immersion  ${\Bpsi}^{-1}: \BLa'\ra \BF$
  of degree bounded by  $d^5$;

\item [(P3$'$)] An immersion $\BUps\ra \BV^n$;

\item[(P4$'$)] An immersion $\hat\BE^{i+1} \ra  \BLa' $  which restricts to
   a homeomorphism $K_{E^{i+1}}\ra   K_\La$.

\end{itemize}

Moreover,

\begin{itemize}
\item[(P5$'$)] The images of the  immersions
 $\BLa_k'\ra \BV$ are disjoint from the Julia pieces $K_{\La_j}$,  $k\not=j$.

\item[(P6$'$)]
There are at least $m/2$ domains $\La_k$.

\end{itemize}
******************************************}
